%% LyX 2.4.4 created this file.  For more info, see https://www.lyx.org/.
%% Do not edit unless you really know what you are doing.
\documentclass[english]{article}
\usepackage[T1]{fontenc}
\usepackage[latin9]{inputenc}
\usepackage{enumitem}
\usepackage{amsmath}
\usepackage{amsthm}
\usepackage{geometry}
\geometry{verbose,tmargin=1in,bmargin=1in,lmargin=1in,rmargin=1in}

\makeatletter
%%%%%%%%%%%%%%%%%%%%%%%%%%%%%% Textclass specific LaTeX commands.
\numberwithin{equation}{section}
\numberwithin{figure}{section}
\newlength{\lyxlabelwidth}      % auxiliary length 

%%%%%%%%%%%%%%%%%%%%%%%%%%%%%% User specified LaTeX commands.
\usepackage{fancyhdr}
\usepackage{lastpage}
\pagestyle{fancy}
\usepackage{enumitem}
\usepackage{ifthen}
\everymath=\expandafter{\the\everymath\displaystyle}
%\DeclareMathSizes{10}{10}{10}{10}
\usepackage{multicol}

\usepackage{tikz}
\usepackage{tikz-3dplot}
\usetikzlibrary{calc,arrows}

\usetikzlibrary{patterns}
\usetikzlibrary{plotmarks}
\usepackage{pgfplots}
\pgfplotsset{compat=newest}

\usepackage{calc}
\usepackage[nomessages]{fp}

\usepackage{datetime}
% Added by lyx2lyx
\setlength{\parskip}{\smallskipamount}
\setlength{\parindent}{0pt}

\makeatother

\usepackage{babel}
\begin{document}
\renewcommand{\author}{Dr.\,James Ashe}

\newcommand{\classNum}{336}

\newcommand{\classSec}{A}

\newcommand{\examType}{HW 5.1}

\renewcommand{\date}{\formatdate{6}{11}{2013}}

%%First page's header%%%%%%%%%%%%%%%%%%%%%%
\fancypagestyle{firststyle} %%header settings for first pagr
{
	\fancyhead[L]{MTH \classNum{}(\classSec{}) \author{}\\
	\textbf{\examType{}} ~\date{}}
	\fancyhead[C]{}
	\fancyhead[R]{\textbf{Solutions}}
	\setlength{\headsep}{14pt}
}
%%%%%%%%%%%%%%%%%%%%%%%%%%%%%%%%%%%%%%%%%%

%%Next page's header%%%%%%%%%%%%%%%%%%%%%%%
\setlist{nolistsep}
\lhead{\classNum{}(\classSec{})} 
\chead{\textbf{Solutions}} 
\cfoot{\thepage{}~of~\pageref{LastPage}} 
\setlength{\headsep}{5pt} 
%%%%%%%%%%%%%%%%%%%%%%%%%%%%%%%%%%%%%%%%%%%

%%Various settings; see the preamble for other settings%%%%%
%\setenumerate[0]{leftmargin=12pt,itemindent=0pt,labelwidth=0pt}
\setlist{topsep=.5pt}
\renewcommand{\labelenumi}{\textbf{\arabic{enumi}.}}
\renewcommand{\labelenumii}{\textbf{\alph{enumii}.}}
\thispagestyle{firststyle}
%%%%%%%%%%%%%%%%%%%%%%%%%%%%%%%%%%%%%%%%%%

\newcommand{\xmax}{0}
\newcommand{\xmin}{-\xmax}
\newcommand{\xstep}{0}
\newcommand{\ymax}{0}
\newcommand{\ymin}{-\ymax}
\newcommand{\ystep}{0} 
\newcommand{\dspace}{\vspace{1cm}}
\newcommand{\xa}{0}
\newcommand{\xb}{0}
\newcommand{\ya}{1}
\newcommand{\yb}{1}
\begin{enumerate}[resume]
\item \vspace{0cm}

\begin{enumerate}[resume]
\item $\det\begin{bmatrix}\begin{array}{rrr}
-1 & 3 & 5\\
6 & 4 & 2\\
-2 & 5 & 1
\end{array}\end{bmatrix}=\det\begin{bmatrix}\begin{array}{rrr}
-1 & 3 & 5\\
0 & 22 & 32\\
0 & -1 & -9
\end{array}\end{bmatrix}=-\det\begin{bmatrix}\begin{array}{rrr}
-1 & 3 & 5\\
0 & -1 & -9\\
0 & 22 & 32
\end{array}\end{bmatrix}=-\det\begin{bmatrix}\begin{array}{rrr}
-1 & 3 & 5\\
0 & -1 & -9\\
0 & 22 & 32
\end{array}\end{bmatrix}=-\det\begin{bmatrix}\begin{array}{rrr}
-1 & 3 & 5\\
0 & -1 & -9\\
0 & 0 & -166
\end{array}\end{bmatrix}=-(-1)(-1)(-166)=166$
\item $\det\begin{bmatrix}\begin{array}{rrrr}
1 & -1 & 0 & 1\\
0 & 2 & 1 & 1\\
2 & -2 & 2 & 3\\
0 & 0 & 6 & 2
\end{array}\end{bmatrix}=\det\begin{bmatrix}\begin{array}{rrrr}
1 & -1 & 0 & 1\\
0 & 2 & 1 & 1\\
0 & 0 & 2 & 1\\
0 & 0 & 6 & 2
\end{array}\end{bmatrix}=\det\begin{bmatrix}\begin{array}{rrrr}
1 & -1 & 0 & 1\\
0 & 2 & 1 & 1\\
0 & 0 & 2 & 1\\
0 & 0 & 0 & -1
\end{array}\end{bmatrix}=(1)(2)(2)(-1)=-4$
\end{enumerate}
\item $A$ is square. So if $A$ has a row of zeros then it is singular
which implies that $\det A=0$ by Theorem 1.2.\\
Alternatively, Suppose that the $j^{\mbox{th}}$ row is all zeros,
$\mathbf{A}_{j}=\mathbf{0}$. Then $2\mathbf{A}_{j}=\mathbf{A}_{j}=\mathbf{0}$.
Now apply property (2) of proposition 1.1 to the $j^{\mbox{th}}$
row so that $\det A=2\det A\Rightarrow\det A=0$.\\
Note that the same result can be shown for a matrix with a column
of zeros using Corollary 1.8. 
\item [\textbf{5.}]$cA$ means that $c$ multiples each of the $n$ rows
of $A$. So then by property (2) of proposition 1.1, $\det(cA)=\underset{n}{\underbrace{c\cdot c\cdots c}}\det A=c^{n}\det A$.
\item [\textbf{9.}]We need to assume the results from parts (b). If $k=\ell,$
then from part (b),
\begin{eqnarray*}
\det\begin{bmatrix}\begin{array}{r|r}
A & B\\
\hline C & D
\end{array}\end{bmatrix} & = & \det A\det(D-CA^{-1}B)\mbox{.}
\end{eqnarray*}
By Theorem 1.5, $\det A\det(D-CA^{-1}B)=\det(AD-ACA^{-1}B)$. It is
given that $CA=AC$ (this assumption is necessary) so then 
\begin{eqnarray*}
\det A\det(D-CA^{-1}B) & = & \det(AD-(AC)A^{-1}B)\\
 & = & \det(AD-(CA)A^{-1}B)\\
 & = & \det(AD-CB)
\end{eqnarray*}
 as needed. 
\end{enumerate}

\end{document}
